% Generated draft: requires mathematical and source review.
\documentclass[11pt]{article}
\usepackage[margin=25mm]{geometry}
\usepackage{fontspec}
\setmainfont{Latin Modern Roman}
\usepackage{amsmath,amssymb,mathtools,unicode-math}
\setmathfont{Latin Modern Math}
\usepackage{enumitem,needspace,xurl,hyperref}
\setlength{\parindent}{0pt}
\setlength{\parskip}{5pt}
\setlength{\emergencystretch}{3em}
\newcommand{\N}{\mathbb{N}}
\newcommand{\Z}{\mathbb{Z}}
\newcommand{\Q}{\mathbb{Q}}
\newcommand{\R}{\mathbb{R}}
\newcommand{\C}{\mathbb{C}}
\newcommand{\F}{\mathbb{F}}
\newcommand{\A}{\mathbb{A}}
\newcommand{\PP}{\mathbb{P}}
\newcommand{\E}{\mathbb{E}}
\newcommand{\eps}{\varepsilon}
\newcommand{\dd}{\,\mathrm{d}}
\begin{document}
\section*{UnsolvedMath ID: 30006642}
\textbf{Source problem code:} catalogue entry
\subsection{Definitions and mathematical statement}
The supplied record contains two intended research tasks, but omits the equations and hypotheses needed to formulate either as a complete theorem.

Let \(N\in\N\) denote the particle number, and let \(\gamma>0\) and \(\ell>0\) denote the parameters appearing in the source. The requested statements are:
\begin{enumerate}
\item Establish dynamical metastability for the mean-field partial differential equation governing the source's weakly interacting diffusions. The referenced PDE, its spatial domain, admissible initial data, interaction assumptions, and a quantitative definition of dynamical metastability are not supplied.
\item Establish convergence of the source's particle system to the massive Arratia flow in the joint regime
\[
N\longrightarrow\infty,\qquad
\gamma\longrightarrow\infty,\qquad
\ell\longrightarrow0,\qquad
\frac{\gamma\ell}{\log N}\longrightarrow\infty.
\]
The intended joint limit can be expressed using sequences \((N_j,\gamma_j,\ell_j)\), with \(N_j\geq2\), satisfying all four limits as \(j\to\infty\). Whether the assertion is intended for every such sequence, and whether additional restrictions are required, must be checked against the complete source equations. The particle dynamics, initial configurations, time and space rescalings, limiting-flow definition, state space, and mode of convergence are absent from the supplied record.
\end{enumerate}

The following source heuristics qualify these tasks; they are not asserted here as proved estimates. After initial clustering, the source describes clusters with particle-index sets \(I_k\), masses
\[
m^{(k)}=\frac{|I_k|}{N},
\]
and centers \(\overline X_t^{(k)}\). It reports an approximate cluster count \(K\approx k_{\max}\), but does not define \(k_{\max}\). Invoking symmetry of an interaction profile \(w\), it gives the approximate center dynamics
\[
\dd\overline X_t^{(k)}
\approx\sqrt{\frac{2}{N m^{(k)}}}\,\dd W_t^{(k)},
\]
with independent Brownian motions and diffusivity \((Nm^{(k)})^{-1}\). Centers are said to merge when their separation becomes comparable to \(\ell\), with \(t_{\mathrm{coalesce}}\sim N\). The source expressly identifies this coalescence as a finite-\(N\) effect invisible to the mean-field PDE. Consequently, this mechanism must not be silently identified with metastability of that PDE.

For a settled \(K\)-cluster state, the source describes slow exchange of mass between neighboring clusters. Its formal Gaussian approximation assigns cluster \(k\) variance
\[
\sigma_k^2=\frac{\ell}{w''(0)m^{(k)}\gamma},
\]
and, with \(\Delta w=|w(0)|\), an effective barrier of height
\[
\gamma\ell\,\Delta w\,m^{(k)}.
\]
The accompanying Eyring--Kramers heuristic gives a mass-escape timescale proportional on an Arrhenius scale to
\[
\exp\!\bigl(\gamma\ell\,\Delta w\,m^{(k)}\bigr).
\]
These expressions require hypotheses not supplied here, including sufficient regularity of \(w\) and positivity of \(w''(0)\) for the stated variance of a positive-mass cluster. The extracted heading separately renders the coarsening scale as \(e^{\gamma\ell\Delta w}/m_{\mathrm{crit}}\); its exponent grouping is uncertain and \(m_{\mathrm{crit}}\) is undefined. No identification of this heading with the mass-dependent escape scale is justified by the record.

Finally, the excerpt describes the particle system as ergodic with respect to a Gibbs measure involving \(W_{\gamma,\ell}\), and labels a microscopic-reversibility scale \(t_{\mathrm{reversibility}}\sim e^{N\Delta F}\). The displayed Gibbs density is corrupted, \(W_{\gamma,\ell}\) and \(\Delta F\) are undefined, and the subsequent qualification concerning a discontinuous phase transition is truncated. These statements supply context, not complete assumptions or established conclusions for this draft.
\subsection{Short English statement}
The intended problems are to prove metastable behavior of the mean-field PDE and a massive Arratia flow limit under the four stated simultaneous limits. The source distinguishes finite-particle cluster coalescence from slow mass exchange. Missing dynamics, hypotheses, rescalings, and convergence definitions prevent a complete theorem statement; the supplied timescales remain heuristics requiring human mathematical review.
\subsection{Sources}
\begin{itemize}
\item Local UnsolvedMath registry, permanent ID \(30006642\), source code \texttt{catalogue entry}, title: \emph{Metastability and Massive-Arratia Limits for Weakly Interacting Diffusions}; collection: Oberwolfach Reports Open Problems; proposed year: 2026. No individual proposer is supplied.
\item Supplied originating reference: \emph{Flows on Measure Spaces and Applications in Machine Learning} (2026), Oberwolfach Reports, DOI 10.4171/owr/2026/13, pp.~795--884. \url{https://doi.org/10.4171/owr/2026/13}. The excerpt refers to particle system (1), mean-field PDE (2), and reference [6], but does not supply those equations or the bibliographic entry for [6].
\item Supplied literature source: Konarovskyi--von Renesse, \emph{Modified massive Arratia flow}. \url{https://arxiv.org/abs/1504.00559}. Its identification with the excerpt's reference [6] is not established by the supplied record.
\item Literature assessment dated exactly \texttt{2026-08-22}: the supplied curation labels the problem open, states that the originating work explicitly presents the joint-scaling convergence as an open question, and characterizes existing Arratia-flow papers as studies of the limit model rather than this particle-system justification. This is supplied dated curation, not an independently verified literature assessment. The background requests independent verification of time-sensitive status and cited claims.
\item The registry supplies the statement-status label \texttt{corrected\_verified} and reports that the statement was reconstructed and checked against the official report. These are supplied provenance assertions; this draft does not independently verify them or the mathematical claims.
\end{itemize}
\end{document}
